% Auto-generated editable source. Edit freely in VS Code.
% This is a parallel source; the live site is still the .md / .html file.
\documentclass[11pt]{article}
\usepackage[margin=1in]{geometry}
\usepackage{graphicx}
\usepackage{hyperref}
\usepackage{enumitem}
\usepackage{xcolor}
\newcommand{\finding}[1]{\par\medskip\noindent\textbf{Finding.} #1\par\medskip}
\newcommand{\stat}[2]{\par\medskip\begin{center}{\LARGE\textbf{#1}}\\[2pt]{\small #2}\end{center}\par\medskip}

\title{Interpreting and evaluating protein language models}
\author{Saanvi S. Subramanian}
\date{}
\begin{document}
\maketitle

\section*{01 - The Challenge}
\section*{A good score is not the same as a good protein}
The central question\par
Proxy
  a computable score used in place of a slow measurement.
  Wet-lab measurements are slow, so a model score stands in for binding, stability, or function. That shortcut is useful, and it also means the score, not the property, is what a search actually optimizes against.\par
Embedding effect versus validated biology
  a shift in a model's representation is a fact about the model, not yet a fact about the protein.
  Much of this work turns on that distinction. An embedding can move in a consistent direction while the move is explained by ordinary sequence properties, so the burden is to show that a residual survives the obvious controls before calling it biology.\par
\section*{02 - The Approach}
\section*{An audit battery, run before the number is believed}
\section*{03 - Results}
\section*{Seven threads, one thesis}
\subsection*{Auditing Protein-Function Prediction}
\textit{No BLAST, New Problems! Auditing Homology-Free Protein-Function Prediction}\par
Finding 1\par
72\%\par
of the embedding movement is reproduced by a composition-matched control\par
92.5\%\par
of the displacement variation is explained by protein length alone\par
Finding 2\par
Why a naturalness metric?\par
5 SD\par
how far an unconstrained optimizer pushes the proxy past the natural-protein floor\par
Finding 3\par
Finding 4\par
\textbf{Limitations.} The audit is demonstrated on one framework and one signature-domain finding. Its generality to other frameworks, signatures, and domains is untested here.\par
\subsection*{Do Benchmarks Predict Better Designs?}
\textit{Does Benchmark Correlation Predict Selection Quality for Protein Design?}\par
Why this metric?\par
0.50 -> 0.16\par
squared rank association with selection value, from the top 10\% down to the top 1\%\par
Finding 1\par
Finding 2\par
[TODO: Saanvi to add Methods appendix, assay selection, scoring procedures, prediction-model specs, supervised-model ingredients, and CV schemes.]\par
\textbf{Limitations.} The comparisons are statistically dependent: 36 assay-scheme comparisons reuse 12 assays, so the bootstrap intervals are likely optimistic. What is next is a theoretical and causal analysis of how correlation maps to selection, and broader corpora and scorer types.\par
\subsection*{Identifiability in Protein-Design Agents}
\textit{Representational Identifiability Limits Action Selection in a Protein-Design Agent}\par
R\textasciicircum{}2 = 0.870\par
of pairwise destination distance explained by a translation placeholder count, in ESM-2 650M\par
Finding 1\par
Finding 2\par
[TODO: Saanvi to add the placeholder-truncation figure; its caption must state the threshold choice explicitly and consistently.]\par
\subsection*{Red-Teaming Agent Verifiers}
\textit{When the Verifier Is Wrong: Red-Teaming Agent Verification Where Ground Truth Exists}\par
Threat model\par
0.991 AUROC, 0.028 gained\par
a fully leaked scorer looks near-perfect yet adds almost no good candidates per batch over random\par
Finding 1\par
Finding 2\par
Finding 3\par
[TODO: Saanvi to add the detection-versus-recovery table across probe sizes with full units, and a Methods appendix with batch size, ridge features and regularization, and the probe rejection rule.] What is next is mapping these attacks onto an LLM-judge setting.\par
\subsection*{Do Model Scores Mean Better Proteins?}
\textit{Do Protein Model Scores Predict Better Proteins? Testing In-Silico Evidence Against Experimental Measurements}\par
[TODO: Saanvi to fill in abstract, key numbers, and figures. OpenReview reviews can be pulled on request.]\par
\subsection*{More Than Composition?}
\textit{Does a Protein Design Pipeline Use More Than Composition?}\par
[TODO: Saanvi to fill in abstract, key numbers, and figures. OpenReview reviews can be pulled on request.]\par
\subsection*{Do Pooled Activations Explain Interventions?}
\textit{Do Pooled Activations Help Explain Intervention Outcomes?}\par
[TODO: Saanvi to fill in abstract, key numbers, and figures. OpenReview reviews can be pulled on request.]\par
\section*{Takeaways}
\section*{What the threads add up to}
Controls before conclusions\par
Correlation is not selection\par
Identifiability has limits\par
Verifiers can be gamed\par
\textbf{Limitations across the line.} Each result is scoped to the framework, model, and search setup it was measured in. The audits are demonstrated rather than proven general, several threads reuse a small set of assays or signatures, and the appendices marked above still need to be filled before any of these numbers should be quoted outside the papers.\par
\section*{Cite}
@inproceedings\{subramanian2026noblast,
  author    = \{Saanvi S. Subramanian\},
  title     = \{No BLAST, New Problems! Auditing Homology-Free Protein-Function Prediction\},
  booktitle = \{Women in Machine Learning Workshop (WiML), NeurIPS\},
  year      = \{2026\},
  url       = \{https://openreview.net/forum?id=hsb1l06LlF\}
\}

@inproceedings\{subramanian2026benchmark,
  author    = \{Saanvi S. Subramanian\},
  title     = \{Does Benchmark Correlation Predict Selection Quality for Protein Design?\},
  booktitle = \{NeurIPS Workshops (Trust-AI-Eval; GEM Bio)\},
  year      = \{2026\},
  url       = \{https://openreview.net/forum?id=KxSFpyRXzp\}
\}

@inproceedings\{subramanian2026identifiability,
  author    = \{Saanvi S. Subramanian\},
  title     = \{Representational Identifiability Limits Action Selection in a Protein-Design Agent\},
  booktitle = \{NeurIPS Agentic Life Sciences Workshop (AgenticLS)\},
  year      = \{2026\},
  url       = \{https://openreview.net/forum?id=GhNJ7TTUyi\}
\}

@inproceedings\{subramanian2026verifier,
  author    = \{Saanvi S. Subramanian\},
  title     = \{When the Verifier Is Wrong: Red-Teaming Agent Verification Where Ground Truth Exists\},
  booktitle = \{NeurIPS Verify-Agents Workshop\},
  year      = \{2026\},
  url       = \{https://openreview.net/forum?id=RyzGAFDM8S\}
\}\par
\section*{Cite}
\begin{verbatim}
@inproceedings{subramanian2026noblast,
  author    = {Saanvi S. Subramanian},
  title     = {No BLAST, New Problems! Auditing Homology-Free Protein-Function Prediction},
  booktitle = {Women in Machine Learning Workshop (WiML), NeurIPS},
  year      = {2026},
  url       = {https://openreview.net/forum?id=hsb1l06LlF}
}

@inproceedings{subramanian2026benchmark,
  author    = {Saanvi S. Subramanian},
  title     = {Does Benchmark Correlation Predict Selection Quality for Protein Design?},
  booktitle = {NeurIPS Workshops (Trust-AI-Eval; GEM Bio)},
  year      = {2026},
  url       = {https://openreview.net/forum?id=KxSFpyRXzp}
}

@inproceedings{subramanian2026identifiability,
  author    = {Saanvi S. Subramanian},
  title     = {Representational Identifiability Limits Action Selection in a Protein-Design Agent},
  booktitle = {NeurIPS Agentic Life Sciences Workshop (AgenticLS)},
  year      = {2026},
  url       = {https://openreview.net/forum?id=GhNJ7TTUyi}
}

@inproceedings{subramanian2026verifier,
  author    = {Saanvi S. Subramanian},
  title     = {When the Verifier Is Wrong: Red-Teaming Agent Verification Where Ground Truth Exists},
  booktitle = {NeurIPS Verify-Agents Workshop},
  year      = {2026},
  url       = {https://openreview.net/forum?id=RyzGAFDM8S}
}
\end{verbatim}
\end{document}